% SPDX-License-Identifier: Apache-2.0
\section{What we got wrong}

The decision log records \ClaimDecisionEntries{} entries. These are the ones that changed a result, listed
because a study that reports only its surviving claims gives a reviewer no way to judge how
hard the others were pushed.

\paragraph{A certificate that appeared to work was vacuous.} An early run certified a large
set of configurations. Every selected $\lambda$ sat exactly at the upper band edge --- where
band membership is $f(x) < \tauhi$ and the in-band rule is $g(x) \ge \lambda$, so the flagged
set is empty by construction and the risk is structurally zero. It was the only grid point
that could clear the Holm level, and it certified nothing. The decision grid now excludes the
upper endpoint, and the body reports the smallest margin below $\tauhi$ (\ClaimCertificateMargin)
so the property is visible rather than assumed.

\paragraph{The same certificate rested on an invalid $p$-value.} The risk was formulated as
$\max(0,\ \text{floor} - \text{recall})/\text{floor}$. Hoeffding--Bentkus bounds a mean of
per-observation losses; that expression is a nonlinear transform of a mean, so the quantities
it produced were not $p$-values and nothing resting on them was a certificate.

\paragraph{The central result was narrowed twice.} A single seed's maximum was first reported
as a central estimate of the temporal inflation. Re-running across seeds
gave a smaller figure with a real interval; re-running across five rolling calibration origins
then showed the deviation is \emph{origin-dependent}, with the earliest origin conservative.
The clean headline did not survive, and the body reports the range instead.

\paragraph{A power gate that could not fail.} The pre-registration requires the minimum
detectable effect \emph{before} the quantum comparison. The calculation had no caller at all,
so H4 was going to be observed rather than tested. The fix estimated noise from the baseline
against itself --- correct reasoning, wrong quantity, because H4 compares model \emph{families}
and their measured standard error (\ClaimMpsStandardErrorBestChi) is \ClaimPowerSeRatio{}-fold
larger than the seed-to-seed figure. The gate passed a comparison it should have flagged. It
now estimates cross-family noise from a second classical family (\ClaimPowerSeCrossFamily{},
giving a pre-registered minimum detectable effect of \ClaimPowerMde{}) and binds on the larger
of the two, and H4 is reported as underpowered against the effect the comparison could
actually resolve.

\paragraph{A claim asserted from theory and withdrawn when measured.} We stated that ignoring
cluster structure makes intervals too narrow by roughly the square root of the cluster size.
That is the classical result for a sample mean. Average precision is a rank statistic; on a
fixture with cluster-level risk the clustered interval came out \emph{narrower} than the
row-level one. The row-level interval is invalid under this
dependence, but the sign of its error is statistic-dependent, and we no longer claim otherwise.

\paragraph{An answer to a named objective, confounded by the base rate.} Section~4 reported the
tensor network at \ClaimMpsShareBand\,\% of the baseline in the band against
\ClaimMpsShareFull\,\% at full scale, and concluded the deficit grew with dimension. It does
not: average precision floors at the base rate, and the band is \ClaimBandFraudRate\,\%
positive against \ClaimTestFraudRate\,\%, so the two ratios start from different floors. As
lift over each floor they agree --- \ClaimMpsLiftShareBand\,\% and \ClaimMpsLiftShareFull\,\%.
Written as the answer to \Cited{secondary objective 4.2}, and withdrawn.

\paragraph{Two silent bugs in the tensor network.} The initialisation placed the identity on
both physical channels, making the contraction input-independent; logits were bit-identical
across samples, so the AUC came out at exactly chance. Separately, closing the final bond to
dimension one made renormalisation map every sample to $\pm 1$. Both produced plausible
near-chance numbers rather than errors. A third latent hazard was removed while diagnosing a
non-finite loss at 431 sites: a discarded log-norm was multiplied by zero, and
$\infty \times 0$ is not zero.
